\section{Versuchsaufbau und -durchführung}
\textcolor{blue}{\textit{einleitende Worte}}\\

\subsection{Versuchsaufbau}
- 1 Umfragenlink (via survio)\\
- jede Anleitung war enthalten\\
- möglichst eine Anweisung pro Seite, damit der viele Text nicht überflogen wird.\\
- alle benötigten Links für die verschiedenen Aufgaben wurden mitgegeben\\
-> Vergleichbarkeit; Shortcuts vermeiden\\
Ablauf:
- Abfrage der moderierenden Varaiblen (Alter, Arbeitsdauer, Expertise)
- Beantwortung von Wissensfragen zum alten System (Variablen die für das neue system nicht benötigt werden)
für beide Systeme je 1x:
    - System-Vorstellung
    - 4 Aufgaben (6 inkl. Teilaufgaben) auf einzelnen Folien\\
    - Umfragebogen mit Schieberegler (7-Punkte-Skala)\\
- Einschätzung, ob die neue Oberfläche in Zukunft genutzt werden wird (5-Punkte-Skala)
- Möglichkeit zu Anmerkungen, Feedback und Diskussionen

\subsection{Beschreibung der Stichprobe}
- 11 Mitarbeiter aus 2 Abteilungen\\
-> 9 - 30+ Jahre Expertise mit dem System allgemein (verschiedene Test-Systeme, DB = Prototyp einer bestimmten Ausprägung) [1x 1-5 Jahre; 1x 21-30 Jahre]\\
-> davon 5 (von den 9) - Expertise speziell in dem abgefragten Bereich vorhanden\\
-> Alter: 7x 60+; 3x 50-59; 1x 30-39\\
-> 8x (fast) tägliche Nutzung; 2x regelmäßig (oft \& mittel), 1x unregelmäßig\\
- arbeiten alle mehr oder weniger mit dem alten System\\
- keiner hat das Dashboard zuvor gesehen\\
-> Kollegen die das Projekt betreut haben, sind aus der Untersuchung rausgefallen\\

\subsection{Versuchsdurchführung}
- per MS Teams-Meeting; Bildschirmübertragung aufgenommen (1:1)\\
- Moderator hat möglichst wenig dazwischen geredet\\
-> Bei Versuch zu viel zu reden, unterbinden\\
-> bei Problemen, Hilfestellung geben\\
- die Gruppe wurde geteilt, um einen Reihenfolgeneffekt \textcolor{red}{\cite{Quelle}} zu vermeiden\\
-> 7 Personen haben die Aufgaben zuerst mit dem alten System bearbeitet
-> 4 Personen haben die Aufgaben zuerst mit dem neuen System bearbeitet
